
\subsection{Detailing Related Work}



\cite{du2020performance} Du, N., Liang, Z., Huang, Y., Guo, Z., Yang, H., Wang, S.: Performance optimi-
sation method of pbft consensus for supply chain integration svm. In: 2020 7th
international conference on dependable systems and their applications (DSA). pp.
371–377. IEEE (2020)

Focuses on addressing the performance bottlenecks caused by malicious or dishonest nodes acting as primary nodes.
By integrating SVM with the consensus mechanism, the system can dynamically evaluate node trustworthiness and 
optimize the selection process to improve throughput and reduce latency









6 Eischer,M.,Distler,T.:Latency-awareleaderselectionforgeo-replicatedbyzantine
fault-tolerant systems. In: 2018 48th Annual IEEE/IFIP International Conference
on Dependable Systems and Networks Workshops (DSN-W). pp. 140–145. IEEE
(2018)

latency-
aware mechanism to select the leader of a geo-replicated BFT
system based on end-to-end response times measured by clients.
To prevent faulty replicas from gaining an unfair advantage by
sending protocol messages early, ARCHER relies on a hash-chain-
based approach that enables clients to detect if a protocol phase
has been skipped. In addition, ARCHER offers means to tolerate
incorrect latency values reported by faulty clients and can also be
extended to solve other selection problems such as the placement
of active and passive replicas in resource-efficient BFT systems.








7. Huang, D., Huang, Y., Yang, Y.: Improved pbft consensus algorithm based on
node evaluation and dynamic management. In: Proceedings of the 2023 4th In-
ternational Conference on Big Data Economy and Information Management. pp.
851–855 (2023)

Firstly, design a mechanism for evaluating nodes
to calculate the reputation score of nodes, and then classify them
according to their reputation scores. Nodes with lower reputation
scores can receive consensus outcomes and are not involved in
the consensus process.
Node reputation evaluation is determined by analyzing their con-
duct, which is influenced by various factors.









9. Qiao, P.: Enhancing manufacturing-as-a-service supply chain transparency with
master-slave multi-chain and sqs-pbft. IEEE Access 12, 186605–186616 (2024)

artigo bem questionavel ...










10. Shan, B., Gao, T., Song, L., Cui, Y.: Grouped byzantine fault-tolerant consensus
mechanism based on node behaviour analysis. In: 2025 4th International Sympo-
sium on Computer Applications and Information Technology (ISCAIT). pp. 1890–
1895. IEEE (2025)

he mechanism constructs a
node voting weight evaluation model by introducing a dynamic
analysis of node behavior, and accumulating behavioral data
based on node performances in the consensus process. Second,
based on the performance values of the nodes, the nodes are
divided into high-performance group (HPG), medium-
performance group (MPG), and low-performance group (LPG),
and the master node is randomly selected from the HPG group.
Finally, according to the node grouping, the PBFT full consensus
process is optimized in stages to reduce the communication
overhead and consensus delay.






 
13. Wen, X., Yang, X.: Mapbft: multilevel adaptive pbft algorithm based on discourse
and reputation models. The Computer Journal 68(6), 635–648 (2025)

MAPBFT initially uses
a reputation model to evaluate node parameters such as past performance, reliability, availability, and response delay, providing
predictive insights for the adaptive algorithm. The adaptive algorithm then employs a multi-layer perceptron to predict the reputation
scores of nodes. This enables the selection of high-reputation nodes for consensus participation, narrowing the consensus scope,
and reducing communication overhead. 









15. Yu, L., Wu, Y., Lu, J., Li, T.: An adaptive reputation update mechanism for
primary nodes in pbft. In: 2024 IEEE 23rd International Conference on Trust,
Security and Privacy in Computing and Communications (TrustCom). pp. 1948–
1953. IEEE (2024)

Nodes with lower ranks are eliminated
and lose the opportunity to become the primary node until
all nodes have been eliminated. Nodes that are not eliminated
continue to have the opportunity to become the primary node.

>>>>>>  nao entendi 
               Committee Eviction:    parece ter 
               Nodes with lower ranks are eliminated
and lose the opportunity to become the primary node until
all nodes have been eliminated. Nodes that are not eliminated
continue to have the opportunity to become the primary node.

               Rotation Policy:   E ??     seria N ?     a troca de lider é como PBFT, nao ?
               
               
               








16. Zhang, G., Pan, F., Tijanic, S., Jacobsen, H.A.: Prestigebft: Revolutionizing view
changes in bft consensus algorithms with reputation mechanisms. In: 2024 IEEE
40th International Conference on Data Engineering (ICDE). pp. 1930–1943. IEEE
(2024)

View changes can be invoked by policy-defined criteria
and failure detection. The former can be implemented dif-
ferently according to application specifications. For exam-
ple, a throughput-threshold policy that changes a view if a
leader fails to operate at an expected throughput (e.g., Aard-
vark [22]) or a timing policy that changes a view every 5
minutes. 








17. Zhao, F., Wang, Y.: Pbft consensus algorithm based on reward and punishment
mechanism. In: 2022 3rd International Conference on Information Science and Ed-
ucation (ICISE-IE). pp. 168–173. IEEE (2022)







18. Zhu, X., Hu, X., Zhu, W.: Rgpbft: A reputation-based pbft algorithm with node
grouping strategy. Arabian Journal for Science and Engineering 50(15), 11837–
11850 (2025)






